\chapter{Large Sample Theory for Moment Estimators}

This chapter discusses the asymptotic distribution of estimators obtained through the Method of Moments (MOM). From a mathematical perspective, the treatment here is straightforward: we simply combine the standard \textbf{Central Limit Theorem (CLT)} and the \textbf{Delta Method}.

\paragraph{Motivation}%
\label{par:Motivation}

Suppose you live in 1870. You are an informed statistician, and you know one thing for sure: \textbf{how to estimate expectations}. The Law of Large Numbers (LLN) guarantees that sample averages converge to true expectations.
However, not every estimation problem presents itself directly as an expectation. The Method of Moments (introduced in Chapter 2) is essentially a ``vehicle" for relating arbitrary estimation tasks \textit{back} to the estimation of expectations. Once we do that, we can leverage the powerful tools of the LLN and CLT.

\section{Estimating an Expectation}

The foundation of the Method of Moments is the estimation of a simple expectation.

\textbf{Setup:}
Let $X_1, X_2, \dots$ be i.i.d. observations from the statistical model:
\[ \mathcal{P} = \{ \mathsf{P}_{\theta} : \theta \in \Theta \subseteq \mathbb{R}^k \}. \]
Let $h$ be a function such that $\mathsf{E}_{\theta}[|h(X)|] < \infty$ for all $\theta \in \Theta$. We consider the estimation of the statistical parameter defined by the expectation of this function:
\[ \gamma(\theta) = \mathsf{E}_{\theta}[h(X)]. \]

\subsubsection{Consistency (Law of Large Numbers)}
We employ the ``plug-in" estimator, which replaces the population expectation with the sample average:
\[ \hat{\gamma}_n := \frac{1}{n} \sum_{i=1}^{n} h(X_i). \]
By the Weak Law of Large Numbers (LLN), this estimator converges in probability to the true parameter:
\[ \hat{\gamma}_n \xrightarrow{p} \gamma(\theta) \quad \forall \theta \in \Theta. \]
\textbf{Terminology:} We say that $\hat{\gamma}_n$ is a \textbf{consistent estimator} of $\gamma(\theta)$. As you collect more data ($n \to \infty$), the estimator converges to the target.

\subsubsection{Asymptotic Normality (Central Limit Theorem)}
To understand the \textit{estimation error}, we look to the Central Limit Theorem. If the variance is finite (i.e., $\mathsf{Var}_{\theta}[h(X)] < \infty$), the CLT implies that the standardized error follows a standard normal distribution asymptotically:
\[ \sqrt{n}(\hat{\gamma}_n - \gamma(\theta)) \xrightarrow{d} \mathcal{N}(0, \mathsf{Var}_{\theta}[h(X)]). \]

This result serves as the building block for more complex Method of Moments estimators, where we will estimate parameters $\theta$ that are functions of these expectations $\gamma(\theta)$.

\section{Method of Moments (General Formulation)}

%\subsection{The General Setup}
While the introductory Method of Moments (\autoref{sec:mom}) typically uses simple powers ($T(x) = x, x^2, \dots$), the general theory allows for any set of chosen ``helper functions."

To estimate a parameter vector $\boldsymbol{\theta} \in \mathbb{R}^k$, we choose $k$ functions $T_1, \dots, T_k$. We define the vector of these statistics as:
\[ T(X) = (T_1(X), \dots, T_k(X))^\top. \]

We then define the \textbf{theoretical expectation map} $\tau$, which links the parameters to the moments:
\[ \tau(\boldsymbol{\theta}) = \mathsf{E}_{\boldsymbol{\theta}}[T(X)] = \begin{pmatrix} \mathsf{E}_{\boldsymbol{\theta}}[T_1(X)] \\ \vdots \\ \mathsf{E}_{\boldsymbol{\theta}}[T_k(X)] \end{pmatrix}. \]

\paragraph{The Estimator Construction:}
The Method of Moments (MOM) operates by matching the theoretical expectations to the empirical data averages. Let $\overline{T}_n$ be the vector of sample means:
\[ \overline{T}_n = \frac{1}{n} \sum_{i=1}^n T(X_i). \]

We obtain the estimator $\hat{\boldsymbol{\theta}}_n$ by solving the system:
\[ \overline{T}_n = \tau(\hat{\boldsymbol{\theta}}_n). \]



If the mapping $\tau$ is injective (one-to-one), we can explicitly define the estimator using the inverse mapping $\tau^{-1}$:
\[ \hat{\boldsymbol{\theta}}_n = \tau^{-1}(\overline{T}_n). \]
\begin{remark}[]
This estimator is well-defined provided the sample average $\overline{T}_n$ falls within the image of $\tau$.
\end{remark}

\subsection*{Asymptotic Distribution (Delta Method Applied to MOM)}

Since the estimator is defined via a differentiable transformation of sample means ($\hat{\boldsymbol{\theta}}_n = \tau^{-1}(\overline{T}_n)$), its asymptotic behavior follows directly from the Delta Method combined with the Central Limit Theorem.

\begin{theo}[Asymptotic Normality of MOM]\label{thm:mom_asymptotics}
Let $X_1, X_2, \dots$ be i.i.d. observations generated by $\mathsf{P}_{\theta_0}$. Assume the following conditions:
\begin{enumerate}[label=\roman*)]
    \item \textbf{Finite Variance:} The chosen statistics have finite second moments, i.e., $\mathsf{E}_{\theta_0}[\|T(X)\|_2^2] < \infty$.
    \item \textbf{Invertibility:} The mapping $\tau$ is injective and continuously differentiable in an open neighborhood of $\theta_0$.
    \item \textbf{Non-Singularity:} The Jacobian matrix $\tau'(\bs{\theta}) = \left( \frac{\partial \tau _{j}(\bs{\theta})}{\partial \theta_{l }}  \right)_{jl}  \in \mathbb{R}^{k \times k}$ is invertible at $\bs{\theta} = \bs{\theta}_0$.
\end{enumerate}

Then, the Method of Moments estimator $\hat{\boldsymbol{\theta}}_n$ exists with probability tending to 1 as $n \to \infty$, and satisfies:
\[ \sqrt{n}(\hat{\boldsymbol{\theta}}_n - \bs{\theta}_0) \stackrel{d}{\longrightarrow} \mathcal{N}_k \Big( 0, \Sigma_{\text{MOM}} \Big), \]
where the asymptotic covariance matrix is:
\[ \Sigma_{\text{MOM}} = \tau'(\bs{\theta}_0)^{-1} \operatorname{Var}_{\bs{\theta}_0}[T(X)] \, \tau'(\bs{\theta}_0)^{-\top}. \]
\end{theo}

\subsubsection{Anatomy of the Variance}
The asymptotic variance formula can be understood as having two distinct components:
\begin{enumerate}
    \item \textbf{Inherited Variance ($\operatorname{Var}_{\theta_0}[T(X)]$):} This is the ``middle" term. It represents the unavoidable error in estimating the expectations of the helper functions $T$. If $T(X)$ has high variance, your input to the equation system is noisy.
    \item \textbf{Error Propagation ($\tau'(\theta_0)^{-1}$):} This term (and its transpose) comes from the Delta Method. It represents the sensitivity of the inverse map $\tau^{-1}$. It answers: ``How much does the solution $\hat{\theta}$ wiggle if the moments $\overline{T}_n$ wiggle?" If the derivative of $\tau$ is small (flat), the inverse derivative is large, amplifying the error.
\end{enumerate}

TODO: improve this proof
\begin{proof}[Application of Inverse Function Theorem]
By the \textbf{Inverse Function Theorem}, if $\tau$ is differentiable with an invertible Jacobian at $\bs{\theta}_0$, then locally around $\tau(\bs{\theta}_0)$, the inverse function $\tau^{-1}$ exists and is differentiable. Crucially, the derivative of the inverse is the inverse of the derivative:
\[ (\tau^{-1})'(t) \Big|_{t=\tau(\bs{\theta}_0)} = [\tau'(\bs{\theta}_0)]^{-1}. \]

By the \textbf{CLT}, the sample moments are asymptotically normal:
   \[ \sqrt{n} (\overline{T}_n - \tau(\theta_0)) \xrightarrow{d} Z \sim \mathcal{N}_k(0, \mathsf{Var}_{\theta_0}[T(X)]). \]

By the \textbf{Delta Method} applied to the function $g = \tau^{-1}$:
   \[ \sqrt{n}(\tau^{-1}(\overline{T}_n) - \tau^{-1}(\tau(\theta_0))) \xrightarrow{d} [\tau'(\theta_0)]^{-1} Z. \]

The covariance of the linear transformation $AZ$ is $A \operatorname{Cov}(Z) A^\top$, yielding the result.
\end{proof}

\section{Application to Exponential Families}

\subsection*{Equivalence of MOM and MLE}

We now consider the specific case where the data follows an \textbf{exponential family}. This setting is particularly elegant because there is a natural connection between the Method of Moments (MOM) and Maximum Likelihood Estimation (MLE).

Consider an exponential family with densities given by:
\[
p_{\bs{\theta}}(x) = \exp\{\langle \bs{\theta}, T(x) \rangle - A(\bs{\theta})\}h(x), \quad \bs{\theta} \in \bs{\theta}.
\]
Here, $\bs{\theta}$ represents the canonical parameters. Suppose the parameter space $\Theta$ is open.



\paragraph{Injectivity of the Mean Map:}
Recall from the properties of exponential families that the cumulant generating function $A(\theta)$ is strictly convex. Consequently, its gradient—which defines the \textbf{mean parametrization} $\tau$—is an injective (one-to-one) mapping:
\[
\tau(\bs{\theta}) := \mathsf{E}_{\bs{\theta}}[T(X)] = \nabla A(\bs{\theta}).
\]
Because this mapping is injective, we can uniquely estimate $\theta$ using the Method of Moments estimator based on the sufficient statistics $T(X)$:
\[
\hat{\bs{\theta}}_{\text{MOM}} = \tau^{-1}(\overline{T}_n).
\]

\begin{claim}
In this setting, the Method of Moments estimator is identical to the Maximum Likelihood Estimator. That is,
\[
\hat{\bs{\theta}}_n = \arg\max_{\bs{\theta} \in \bs{\theta}} \ell_n(\bs{\theta}|\bs{x}).
\]
\end{claim}

\begin{proof}
Since $\Theta$ is open, the MLE (if it exists) must satisfy the first-order condition $\dot{\ell}_n(\theta|\bs{x}) = 0$.
The log-likelihood function for $n$ i.i.d. observations is:
\begin{align*}
\ell_n(\boldsymbol{\theta}|\bs{X}) &= \sum_{i=1}^n \log p_\theta(X_i) \\
&= \sum_{i=1}^n \left[ \langle \boldsymbol{\theta}, T(X_i) \rangle - A(\boldsymbol{\theta}) + \log h(X_i) \right] \\
&= n \langle \boldsymbol{\theta}, \overline{T}_n \rangle - n A(\boldsymbol{\theta}) + \text{const.}
\end{align*}
\textit{Intuition:} The logarithm ``eats" the exponential, exposing the linear structure of the sufficient statistics and the cumulant function.
The score function (gradient) is:
\[
\nabla_{\theta} \ell_n(\boldsymbol{\theta}|\bs{X}) = n \overline{T}_n - n \nabla A(\boldsymbol{\theta}).
\]
Recalling that $\nabla A(\theta) = \tau(\theta)$, the condition for the gradient to be zero is:
\[
n \left[ \overline{T}_n - \tau(\boldsymbol{\theta}) \right] = 0 \iff \overline{T}_n = \tau(\boldsymbol{\theta}).
\]
Thus, solving for the MLE is mathematically equivalent to solving the Method of Moments equation.
\end{proof}

\subsubsection{Fisher Information Gives Asymptotic Variance}

Since we have established that $\hat{\theta}_{\text{MLE}} = \hat{\theta}_{\text{MOM}}$, we can apply the Delta Method results from the previous section to derive the asymptotic distribution of the MLE.

\begin{corollary}[Asymptotic Normality of MLE]
In the considered exponential family setting, the MLE $\hat{\bs{\theta}}_n$ exists with probability tending to 1 as $n \to \infty$, and satisfies:
\[
\sqrt{n}(\hat{\bs{\theta}}_n - \bs{\theta}_0) \stackrel{d}{\longrightarrow} \mathcal{N}_k(0, I(\bs{\theta}_0)^{-1}),
\]
where $I(\bs{\theta}_0)$ is the Fisher Information matrix for a single observation.
\end{corollary}

\subsubsection{Derivation: The ``Sandwich" Cancellation}
The general Method of Moments variance (from \autoref{thm:mom_asymptotics}) takes the form of a ``sandwich":
\[
\text{Asymp. Var} = \underbrace{\tau'(\theta_0)^{-1}}_{\text{Bread}} \underbrace{\operatorname{Var}_{\theta_0}[T(X)]}_{\text{Meat}} \underbrace{\tau'(\theta_0)^{-\top}}_{\text{Bread}}.
\]
In the exponential family context, a miraculous simplification occurs:
\begin{enumerate}[]
    \item 
 \textbf{The Meat (Variance):} The variance of the sufficient statistic is the Hessian of the cumulant function:
    \[ \operatorname{Var}_{\bs{\theta}_0}[T(X)] = \nabla^2 A(\bs{\theta}_0) = I(\bs{\theta}_0). \]
\item   \textbf{The Bread (Jacobian):} The mapping $\tau(\bs{\theta})$ is the gradient $\nabla A(\bs{\theta})$. Therefore, its Jacobian $\tau'(\bs{\theta})$ is the Hessian $\nabla^2 A(\bs{\theta})$:
    \[ \tau'(\bs{\theta}_0) = \nabla^2 A(\bs{\theta}_0) = I(\bs{\theta}_0). \]
\end{enumerate}

Substituting these into the sandwich formula:
\[
I(\bs{\theta}_0)^{-1} \cdot I(\bs{\theta}_0) \cdot I(\bs{\theta}_0)^{-1} = I(\bs{\theta}_0)^{-1}.
\]
This confirms that the MLE achieves the \textbf{Cramér--Rao bound} asymptotically. While the estimator may have finite-sample bias, it reaches the optimal variance level as $n \to \infty$.

\subsubsection{From Theory to Practice: Software and Approximation}

The theoretical result $\sqrt{n}(\hat{\bs{\theta}} - \bs{\theta}_0) \to \mathcal{N}(0, I(\bs{\theta}_0)^{-1})$ motivates the approximation:
\[
\hat{\bs{\theta}} \stackrel{\cdot}{\sim} \mathcal{N}\left(\bs{\theta}_0, \frac{1}{n} I(\bs{\theta}_0)^{-1}\right).
\]
However, this is not directly usable because the true parameter $\bs{\theta}_0$ (and thus the true variance) is unknown.

\paragraph{Implementation Strategy (Slutsky's Theorem):}
In practice (and in statistical software), we replace the unknown information matrix $I(\bs{\theta}_0)$ with a consistent estimator, typically the observed information at the estimated value, $\hat{I} = I(\hat{\bs{\theta}}_n)$. By Slutsky's Theorem, the convergence still holds:
\[
\sqrt{n}\,\hat{\bs{I}}(\hat{\boldsymbol{\theta}}_n)^{1/2}(\hat{\boldsymbol{\theta}}_n-\boldsymbol{\theta}_0) \stackrel{d}{\longrightarrow} \mathcal{N}_k(0,I_k).
\]
This makes MLE extremely ``software-friendly." An optimizer finds the peak of the log-likelihood ($\hat{\bs{\theta}}$), and the Hessian at that peak automatically provides the estimated Fisher Information, which yields standard errors and confidence intervals.

\begin{note}[on Dependent Data]
    
The standard derivation assumes i.i.d. data, where the total information is $I_n(\bs{\theta}) = n I_1(\bs{\theta})$.
For dependent data (e.g., Markov Chains), the Fisher Information does not necessarily scale linearly with $n$. However, generalized theorems exist (see Billingsley, 1961) showing that $\hat{\bs{\theta}} \stackrel{\cdot}{\sim} \mathcal{N}(\bs{\theta}_0, I_{\text{total}}^{-1})$, provided one calculates the total information inherent in the dependent sequence correctly.
\end{note}

\paragraph{Invariance:}
If the model uses a different parametrization $\bs{\lambda}$ (via a diffeomorphism $\bs{\theta}(\bs{\lambda})$), the asymptotic normality is preserved. The asymptotic variance becomes $I(\bs{\lambda}_0)^{-1}$, which can be calculated using the Delta Method and the appropriate Jacobians.
